\section{Intel VTune} 
\label{section:localisation:vtune}

\index{Profiler}
\index{Profiler!Instrumentation}
Intel gives its VTune profiler away free of charge.
The tool provides a fancy GUI that you can use to study your code interactively,
but there's a button at the start screen that tells you how to invoke the
profiler on the command line.
Through this, you can run the profiler on a compute node without the GUI, and
then copy the results back to your workstation for a posteriori assessment.
Once available (on your supercomputer you'll have to load the VTune module,
while a \texttt{source /myinteldir/setvar.sh} sets all variables on a local
installation), a simple \texttt{vtune-gui} brings up the tools.

VTune's various data collection modes are based upon sampling, and there's the
opportunity to dive very deeply into the machine's behaviour.
As the tool displays the code analysis over time and cores, you can zoom into
certain runtime phases and cores manually. 
The delivered data is more detailed on Intel machines, but the tool also is very
useful on AMD.
Some of the more modi require certain system priviliges.


\paragraph*{Overview}
 
\begin{itemize}[noitemsep]
  \item[$\boxplus$] Works also on non-Intel machinery.
  \item[$\boxplus$] Very detailed node- and core-level data including
  performance counter data.
  \item[$\boxplus$] Can trace MPI.
\end{itemize}
 
 
\begin{itemize}[noitemsep]
  \item[$\boxminus$] Some in-depth features allow special system rights that
  normal users typically don't have.
  \item[$\boxminus$] Very easy to get lost in massive amounts of data.
\end{itemize}

